\par\noindent\textbf{(1)} 在平面外布置三个镜像：\(-q\) 于 \((-a,b,0)\)、\(-q\) 于 \((a,-b,0)\)、\(+q\) 于 \((-a,-b,0)\)。加上真实 \(q\) 的四项Coulomb势 \(\Phi(\boldsymbol r)=\sum_jq_j/(4\pi\epsilon_0|\boldsymbol r-\boldsymbol r_j|)\)，在 \(x=0\) 与 \(y=0\) 上各对镜像贡献相互抵消，满足双接地Dirichlet边界条件。
\par\noindent\textbf{(2)} 左侧镜像引力是 \(-q^2/(16\pi\epsilon_0a^2)\hat x\)，下侧镜像引力是 \(-q^2/(16\pi\epsilon_0b^2)\hat y\)。斜对角同号镜像的力为 \(q^2(a\hat x+b\hat y)/[16\pi\epsilon_0(a^2+b^2)^{3/2}]\)。三者叠加给 \(F_x=C[-1/a^2+a/(a^2+b^2)^{3/2}]\)、\(F_y=C[-1/b^2+b/(a^2+b^2)^{3/2}]\)，\(C=q^2/(16\pi\epsilon_0)\)。
\par\noindent\textbf{(3)} 感应电荷随位置变化，故静电机械势能必须取一半像电势：\(U=\tfrac12q\Phi_{\rm images}(a,b)=-C[1/a+1/b-1/\sqrt{a^2+b^2}]\)。分别求 \(-\partial_aU\) 与 \(-\partial_bU\)，恰得到逐镜像Coulomb矢量叠加的两个分量。忽略因子 \(1/2\) 将把机械力错误地放大。
\par\noindent\textbf{(4)} 当 \(a=b\) 时 \(F_x=F_y<0\)，说明电荷趋近两面的交点。令 \(b\to\infty\)，剩余 \(F_x\to-q^2/(16\pi\epsilon_0a^2)\) 且 \(F_y\to0\)，严格恢复单接地平面的镜像力。只有直角等特殊楔角才可用有限数目镜像电荷直接闭合。
